\section{Security}\label{sec:security}

Secrecy until the deadline rests on two assumptions about SHA-256 and two about the attacker's hardware. This section states them, proves the delay in the random-oracle model, and estimates the one hardware parameter that the proof cannot supply.

\subsection{Assumptions}\label{sec:assumptions}

\begin{description}[leftmargin=1.5em,font=\normalfont\bfseries]
  \item[A1. No shortcut.] The chain step and each labelled hash behave like independent random oracles on their whole inputs. We instantiate them heuristically with SHA-256; a chain step is a single compression. If this fails, the delay is lost and the hash must be replaced, which the versioned format allows.
  \item[A2. Speed.] During the delay, no party computes one chain step in less than $\ell$. An error here shrinks the delay in proportion.
  \item[A3. Budget.] The attacker computes at most $R$ hashes per second in total. An error here multiplies the bound of \cref{thm:delay} by the same factor; for the reference puzzle, any factor up to $2^{52}$ keeps it below $2^{-64}$.
  \item[A4. Encryption.] HMAC-SHA256 is a pseudorandom function (\cref{def:prf}).
\end{description}

\subsection{The delay theorem}\label{sec:theorem}

\paragraph{Model.}
The chain steps $F_1, \dots, F_n$ and the labelled hashes of \cref{alg:build,alg:open} are independent random oracles (\cref{sec:rom}). The attacker is a classical randomized algorithm $\Adv = (\Adv_0, \Adv_1)$ with unbounded computation. $\Adv_0$ runs before the puzzle exists and makes $\Qpre$ queries; its whole state passes to $\Adv_1$, which receives the puzzle and makes $\Qon$ queries. Let $Q = \Qpre + \Qon$. Time after publication is counted in oracle rounds: all queries in a round are chosen before any answer from that round arrives.

\paragraph{Honest schedule.}
Computing each chain's mask and accumulator in the same round, $\Open$ derives $K$ in oracle round
\[
  \Tkey = n(m + 2) + 1
\]
and returns it, checked against $C$, one round later. $\Build$ takes $m + n + 3$ rounds: one for the starts, $m$ for the chains, $n$ for the accumulator, with the masks alongside, and one each for $K$ and $C$. Let $P = n(m + 3) + 1$, the number of oracle answers $\Open$ uses, $C$ included.

\begin{theorem}[Classical attackers, random-oracle model]\label{thm:delay}
Let $(Z, K) \gets \Build(n, m, b)$, with randomness independent of the oracles. For every attacker $\Adv = (\Adv_0, \Adv_1)$ as above,
\[
  \Pr\bigl[\Adv_1 \text{ outputs } K \text{ within fewer than } \Tkey \text{ oracle rounds}\bigr]
  \;\le\; \frac{\Qon}{2^{8b}} + \frac{\Qpre}{2^{128 + 8b}} + \frac{QP + P^2 + 1}{2^{256}},
\]
and the first term is absent when $n = 1$.
\end{theorem}

\paragraph{The bound in seconds.}
Assumptions A2 and A3 turn oracle rounds and queries into time. Write $T = n m \ell$, and let the attacker make $\Qon = R \cdot T$ queries after publication. Against $R = 10^{21}$ queries per second, about the hash rate of the whole Bitcoin network (\cref{sec:pricing}), the bound is about $2^{-117}$ for a 30-second lock and $2^{-97}$ for a one-year lock. Precomputation changes this only slightly: a full year of queries at the same rate before a 30-second puzzle is published adds about $2^{-128}$. With 24-byte links the first term, the chance of guessing a link, dominates; for a 30-second lock the full-width terms are smaller by a factor of about $2^{31}$. At a fixed delay the bound depends on the number of chains only weakly, through $P \approx N$, except that the first term is absent when $n = 1$. The bound covers every speed-up, however small: except with this probability, the attacker does not derive the key even one oracle round before the honest opener.

\paragraph{Scope.}
The threshold is the round in which $K$ is derived, not the round in which $\Open$ returns: an attacker that skips the check against $C$ finishes one round before $\Open$ does, with certainty. The precomputation term is needed even when $n = 1$, where the only link is public: an attacker that guesses the salt and the link in advance and computes the key opens the puzzle in zero rounds with probability $2^{-(128 + 8b)}$. The theorem concerns recovery of the key; \cref{thm:hiding} below extends it to the contents. The bounds are classical; extending them to this instantiation against quantum attackers is outside the scope of this paper.


\Cref{app:proof} gives the proof. It follows the structure of Pietrzak and Reyzin's soundness proof for proofs of sequential work~\cite{PR26}.

\paragraph{Hiding the contents.}
An attacker in an encrypted mempool also sees the envelopes sealed under $K$, and may only need to tell two candidate transactions apart. The labelled hashes $\Henc$ and $\Hmk$ that derive an envelope's subkeys are random oracles like the others, but HMAC is computed outside the oracles. This theorem, unlike \cref{thm:delay}, therefore needs the attacker's other computation to be bounded, since an unbounded attacker could evaluate HMAC under every key. We let $\Adv$ make at most $\Qpre$ and $\Qon$ oracle queries, as before, and at most $t$ further computation steps.

\begin{figure}[t]
\centering\small
\newcommand{\gamebox}[2]{\fbox{\begin{minipage}[t][3.1cm][t]{#1}#2\end{minipage}}}
\gamebox{0.27\textwidth}{$\underline{\mathsf{Hide}_\beta(\Adv)}$
\begin{algorithmic}[1]
  \State $\sigma \gets \Adv_0^{H}$
  \State $(Z, K) \gets \Build(n, m, b)$
  \State $i \gets 0$
  \State $\beta' \gets \Adv_1^{H, \mathsf{Enc}}(\sigma, Z)$
  \State \Return $\beta'$
\end{algorithmic}}\hfill
\gamebox{0.33\textwidth}{$\underline{\mathsf{Enc}(M_0, M_1)}$
\begin{algorithmic}[1]
  \State \algorithmicif\ $i = \Qenc$ \algorithmicthen\ \Return $\bot$
  \State \algorithmicif\ $|M_0| \neq |M_1|$ \algorithmicthen\ \Return $\bot$
  \State $i \gets i + 1$
  \State $\eta \gets$ 16 random bytes
  \State \Return $\mathsf{Seal}(K, \eta, M_\beta)$
\end{algorithmic}}\hfill
\gamebox{0.31\textwidth}{$\underline{\mathsf{PRF}_\beta(\mathcal{B})}$
\begin{algorithmic}[1]
  \For{$j = 1, \dots, u$}
    \State $k_j \gets$ 32 random bytes
    \State $g_j \gets \HMAC(k_j, \cdot)$
    \State \algorithmicif\ $\beta = 0$ \algorithmicthen\ $g_j \xleftarrow{\$} \mathcal{F}$
  \EndFor
  \State \Return $\mathcal{B}^{g_1, \dots, g_u}$
\end{algorithmic}}
\caption{The games of \cref{def:hiding,def:prf}. $H$ denotes all the oracles, and $\mathsf{Seal}(K, \eta, M)$ is the envelope of $M$ under $K$ with nonce $\eta$ (\cref{sec:costs}). $\mathcal{F}$ is the set of functions from byte strings to 32-byte values, and $\xleftarrow{\$}$ samples uniformly.}
\label{fig:games}
\end{figure}

\begin{definition}[Encryption game]\label{def:hiding}
For $\beta \in \{0, 1\}$, the game $\mathsf{Hide}_\beta$ and its envelope oracle $\mathsf{Enc}$ are given in \cref{fig:games}. $\Adv_0$ makes at most $\Qpre$ oracle queries; $\Adv_1$ makes at most $\Qon$ and stops within fewer than $\Tkey$ oracle rounds. The advantage of $\Adv$ is
\[
  \Advhide(\Adv) = \bigl|\Pr[\mathsf{Hide}_1(\Adv) = 1] - \Pr[\mathsf{Hide}_0(\Adv) = 1]\bigr|.
\]
\end{definition}

\begin{definition}[Pseudorandomness of HMAC]\label{def:prf}
For $\beta \in \{0, 1\}$ and $u$ keys, the game $\mathsf{PRF}_\beta$ is given in \cref{fig:games}. The advantage of an algorithm $\mathcal{B}$ is
\[
  \Advprf_u(\mathcal{B}) = \bigl|\Pr[\mathsf{PRF}_1(\mathcal{B}) = 1] - \Pr[\mathsf{PRF}_0(\mathcal{B}) = 1]\bigr|,
\]
and $\Advprf_u(t')$ is the largest advantage of an algorithm running in time $t'$.
\end{definition}

\begin{theorem}[Hiding]\label{thm:hiding}
For every attacker as above,
\[
  \Advhide(\Adv)
  \;\le\; 2\Bigl(\frac{\Qon}{2^{8b}} + \frac{\Qpre}{2^{128 + 8b}} + \frac{QP' + P'^2}{2^{256}} + \frac{\Qenc^2}{2^{129}} + \Advprf_{2\Qenc}\bigl(t + O(Q + P')\bigr)\Bigr),
\]
where $P' = P + 2\Qenc$, and the first term is absent when $n = 1$.
\end{theorem}

The first three terms come from \cref{thm:delay}, with $P'$ path values in place of $P$: each envelope adds two subkeys derived from $K$. The fourth is the chance that two nonces repeat, and the last is the cost of replacing HMAC under the $2\Qenc$ subkeys by random functions. For the reference puzzle and few envelopes, the doubled link-guessing term is about $2^{-116}$; the nonce term can dominate when many envelopes share a puzzle. The theorem says nothing from round $\Tkey$ on, when the key can be derived and every envelope under it opens. \Cref{app:hiding} gives the proof.

\subsection{The hardware assumption}\label{sec:pricing}

The attacker's total hash rate can be estimated from deployed hardware: the Bitcoin network computes about $10^{21}$ hashes per second (October 2026), and we take $R = 10^{21}$. A query is one compression, while a Bitcoin hash costs two, since it hashes twice and the first block of the header is computed once; counted in compressions, the network makes about $2 \cdot 10^{21}$ queries per second, which costs one bit: $2^{-116}$ instead of $2^{-117}$ for a 30-second lock. Mining hardware computes many independent hashes at once, of block headers rather than salted chains, and says little about how fast one hash can follow another, so it bounds $R$, not $\ell$.

The speed $\ell$ is the assumption that most needs hardware evidence. \Cref{tab:pricing} collects what is known. Relative to the fastest published design that we know of, a synthesized 7\,nm core at 13.0\,ns per hash~\cite{BCF+20}, we assume an \emph{advantage} of $A \approx 3.7$, that is, $\ell = 3.5$\,ns: 6.6 times faster than one core of the M4 laptop we measured. Some advantage over the published core is certain, since it is a general-purpose design that registers every round and does not exploit the constant part of a chain step's input; \cref{app:hardware} gives a circuit model for a chip built for chains. We do not claim 3.5\,ns is a lower bound. The theorem holds for any value of $\ell$ the reader adopts, and every cost in \cref{sec:costs} scales linearly with $1/\ell$.

\begin{table}[t]
  \centering\small
  \caption{Time per chained SHA-256 hash, from measurements to models.}
  \label{tab:pricing}
  \begin{tabularx}{\textwidth}{@{}>{\raggedright\arraybackslash}p{4.6cm}r>{\raggedright\arraybackslash}X@{}}
    \toprule
    \colhead{Source} & \colhead{Time per hash} & \colhead{Basis} \\
    \midrule
    Apple M4, one core & 23.2\,ns & \tagm{measured}: \leanseal benchmark \\
    AMD EPYC 9554P (Zen 4), one core & 37.4\,ns & \tagm{measured}: \leanseal benchmark, SHA-NI \\
    Published SHA-256 core, 7\,nm~\cite{BCF+20} & 13.0\,ns & \tagm{synthesis}: 67 cycles at 5.15\,GHz; not a fabricated chip \\
    Custom round circuit, 3--5\,nm & 2.9--9.7\,ns & \tagm{model}: \cref{app:hardware}; no such chip has been built \\
    \midrule
    \textbf{Assumed attacker} & \textbf{3.5\,ns} & the design point of this paper ($A \approx 3.7$) \\
    \bottomrule
  \end{tabularx}
\end{table}

\paragraph{Sensitivity.}
Only A1 fails abruptly. An attacker $k$ times faster than assumed opens a 30-second puzzle in $30/k$ seconds; restoring the margin costs $k$ times more sender work and no extra bytes. An attacker with $k$ times the budget multiplies the bound by $k$, which for the reference puzzle stays below $2^{-64}$ unless $k$ exceeds $2^{52}$. Slower honest hardware costs time, never secrecy. For the reference setting the margin is a factor of about 2.5: a 30-second puzzle covers a 12-second window only while one step takes more than $12\,\mathrm{s}/2^{33} \approx 1.4$\,ns. At 1.3\,ns, the advantage $A = 10$ that \cref{app:hardware} does not rule out, it opens in 11.2\,s, before the deadline. The bound of \cref{thm:delay} is negligible next to this; the hardware estimate is what limits the reference parameters. The secrecy is also temporary: a sealed transaction becomes public in any case, so only hardware that exists during its delay matters, and $\ell$ can be repriced as hardware improves.

\subsection{Attacks and link width}\label{sec:attacks}

\begin{itemize}[leftmargin=1.5em]
  \item \textbf{Link guessing.} An attacker can guess a hidden link $x_k$ and run chains $k, k+1, \dots$ while its own earlier chains are still running. This is the first term of the theorem; with 24-byte links each guess is right with probability $2^{-192}$. Because the key needs every chain's end, a guess cannot be checked against $C$ until the attacker's own chains reach link $k$.
  \item \textbf{Precomputation.} Every hash on the honest path depends on the puzzle's fresh salt and a fresh link, so work done before publication or on another puzzle does not help, except as counted by the theorem's second and third terms.
  \item \textbf{Quantum attackers.} Afshar et al.~\cite{ACHLM23} prove that the idealized chained puzzle keeps its delay against quantum openers in the quantum random-oracle model. Carrying this over to \leanseal's exact encoding has not been done; our bound is classical.
  \item \textbf{Early access.} The attacker's clock starts when it first holds the puzzle. Puzzles built in advance must stay private until use.
  \item \textbf{Message substitution.} A claimed key is checked against $C$ with one hash and the tag then authenticates the transaction. But once $K$ is known, anyone can seal a different transaction under the same puzzle, so an ordering must bind the hash of the whole sealed transaction (\cref{sec:flow}).
\end{itemize}

The link width $b$ enters the online guessing term $\Qon / 2^{8b}$ and the preprocessing term. With $n = 16$ and the attacker budget above, 16-byte links (324-byte puzzle) give $2^{-53}$ for a 30-second lock, 24-byte links (452 bytes) give $2^{-117}$, and 32-byte links (580 bytes) give $2^{-148}$; for a one-year lock, which needs at least 33 chains to respect the format's limit $m \le 2^{48}$, the three give about $2^{-33}$, $2^{-97}$ and $2^{-108}$. We use 24 bytes, which keeps the bound near $2^{-96}$ or below for delays up to a year. With $n = 1$ there is no hidden link, so the online guessing term vanishes; the preprocessing term still depends on the salt and link entropy.
